\documentclass[dvisvgm,tikz,border=3pt]{standalone}
\usepackage{amsmath,amssymb}
\usepackage{pgfplots}
\usepackage{CJKutf8}
\pgfplotsset{compat=1.18}
\usetikzlibrary{arrows.meta,calc,positioning}

\definecolor{themebg}{HTML}{2e362c}
\definecolor{themefg}{HTML}{edf4e8}
\definecolor{themegray}{HTML}{9ea897}
\definecolor{themecurve}{HTML}{7eb6ff}
\definecolor{themealert}{HTML}{ff7b7b}
\definecolor{themeamber}{HTML}{f5b942}

\pgfdeclarelayer{background}
\pgfdeclarelayer{foreground}
\pgfsetlayers{background,main,foreground}

\begin{document}
\begin{CJK*}{UTF8}{gbsn}
\pagecolor{themebg}
\begin{tikzpicture}[color=themefg, text=themefg, >=Stealth]

% ── 左子图：主轴旋转与消去交叉项 ──
\begin{scope}[shift={(0,0)}]
  \draw[themegray!70, line width=0.8pt, rounded corners=5pt] (-3.0,-2.6) rectangle (3.0,2.8);
  \node[font=\bfseries\small, text=themecurve, above] at (0, 2.9) {主轴旋转：正交消去交叉项};

  \coordinate (O1) at (0,0);

  \begin{pgfonlayer}{background}
    % 原坐标轴 x1, x2 (细中性灰)
    \draw[->, themegray, line width=0.8pt] (-2.6, 0) -- (2.6, 0) node[right, font=\small] {$x_1$};
    \draw[->, themegray, line width=0.8pt] (0, -2.2) -- (0, 2.4) node[above, font=\small] {$x_2$};

    % 旋转主轴 y1, y2 (旋转30度，暖金实线)
    \draw[->, themeamber, line width=1.4pt] (-2.5*0.866, -2.5*0.5) -- (2.7*0.866, 2.7*0.5) node[above right, font=\bfseries\small] {$y_1$ 主轴};
    \draw[->, themeamber, line width=1.4pt] (2.2*0.5, -2.2*0.866) -- (-2.4*0.5, 2.4*0.866) node[above left, font=\bfseries\small] {$y_2$ 主轴};

    % 倾斜椭圆等值线 (长半轴 2.2, 短半轴 1.1, 倾角 30度)
    \draw[themecurve, line width=2.4pt, rotate=30] (0,0) ellipse (2.2 and 1.1);
    \draw[themecurve!50!themegray, line width=1.0pt, dashed, rotate=30] (0,0) ellipse (1.4 and 0.7);
  \end{pgfonlayer}

  \begin{pgfonlayer}{main}
    \fill[themefg] (O1) circle (2.0pt);
    
    % 半轴标注点
    \coordinate (A1) at ({2.2*cos(30)}, {2.2*sin(30)});
    \coordinate (A2) at ({-1.1*sin(30)}, {1.1*cos(30)});
    \fill[themeamber] (A1) circle (2.0pt);
    \fill[themeamber] (A2) circle (2.0pt);
  \end{pgfonlayer}

  \begin{pgfonlayer}{foreground}
    \node[text=themecurve, font=\bfseries\footnotesize, fill=themebg, inner sep=0.8pt] at (0.2, 1.8) {原方程：$a x_1^2 + 2b x_1 x_2 + c x_2^2 = 1$};
    \node[text=themeamber, font=\bfseries\footnotesize, fill=themebg, inner sep=0.8pt] at (0.2, -1.8) {标准形：$\lambda_1 y_1^2 + \lambda_2 y_2^2 = 1$};
    
    \node[text=themeamber, font=\scriptsize, below right] at (A1) {$\frac{1}{\sqrt{\lambda_1}}$};
    \node[text=themeamber, font=\scriptsize, above left] at (A2) {$\frac{1}{\sqrt{\lambda_2}}$};
  \end{pgfonlayer}
\end{scope}

% ── 右子图：符号惯性与几何形态分类 ──
\begin{scope}[shift={(8.6,0)}]
  \draw[themegray!70, line width=0.8pt, rounded corners=5pt] (-3.4,-2.6) rectangle (3.4,2.8);
  \node[font=\bfseries\small, text=themecurve, above] at (0, 2.9) {符号惯性：定号 vs 不定几何};

  % 上半区：正定 (闭合椭圆)
  \begin{scope}[shift={(-1.2, 1.1)}]
    \draw[themecurve, line width=1.8pt] (0,0) ellipse (1.4 and 0.7);
    \node[text=themecurve, font=\bfseries\small, anchor=west] at (1.6, 0.15) {正定 $(+,+)$};
    \node[text=themegray, font=\scriptsize, anchor=west] at (1.6, -0.25) {闭合椭圆，$\min > 0$};
  \end{scope}

  % 中间区：不定 (双曲线分岔)
  \begin{scope}[shift={(-1.2, -0.3)}]
    \draw[themealert, line width=1.8pt] (-1.3, 0.5) to[bend right=30] (-1.3, -0.5);
    \draw[themealert, line width=1.8pt] (1.3, 0.5) to[bend left=30] (1.3, -0.5);
    \draw[dashed, themegray, line width=0.7pt] (-1.4,-0.5) -- (1.4,0.5);
    \draw[dashed, themegray, line width=0.7pt] (-1.4,0.5) -- (1.4,-0.5);
    
    \node[text=themealert, font=\bfseries\small, anchor=west] at (1.6, 0.15) {不定 $(+,-)$};
    \node[text=themegray, font=\scriptsize, anchor=west] at (1.6, -0.25) {双曲线渐近线零锥};
  \end{scope}

  % 下半区：半正定 (平行退化带)
  \begin{scope}[shift={(-1.2, -1.7)}]
    \draw[themeamber, line width=1.8pt] (-1.4, 0.35) -- (1.4, 0.35);
    \draw[themeamber, line width=1.8pt] (-1.4, -0.35) -- (1.4, -0.35);
    \draw[dashed, themegray, line width=0.7pt] (-1.4, 0) -- (1.4, 0);

    \node[text=themeamber, font=\bfseries\small, anchor=west] at (1.6, 0.15) {半正定 $(+,0)$};
    \node[text=themegray, font=\scriptsize, anchor=west] at (1.6, -0.25) {平行柱面，$\ker A$ 谷底};
  \end{scope}
\end{scope}

% ── 底部总结横条 ──
\node[font=\bfseries\small, color=themecurve, anchor=north] at (4.3, -3.0) {
  核心心智模型：交叉项是坐标系偏离主轴的投影伪像；正交变换保轴长（谱），一般合同保正负（惯性指数）
};

\end{tikzpicture}
\end{CJK*}
\end{document}
